\chapter{Perception}
\label{ch:perception}
\section{Pipeline de perception}
\label{sec:perception-pipeline}

La perception est le premier étage du pipeline : sa tâche est de transformer les images brutes des caméras en une information exploitable par la planification --- avant tout, la \emph{position tridimensionnelle} de l'objet cible dans le repère de base du robot. Elle s'exprime entièrement dans les types de données du chapitre~\ref{ch:architecture} : des \texttt{Frame} en entrée, une \texttt{Scene} peuplée d'\texttt{ObjectInstance} en sortie.

Le traitement enchaîne quatre opérations. D'abord la \emph{capture} : les trois caméras sont lues de façon synchronisée, chacune accompagnée de ses paramètres de calibration et de sa pose dans le repère de base (section~\ref{sec:calibration-cam}). Ensuite la \emph{détection 2D} : dans chaque image, un détecteur localise les objets connus et renvoie, pour chacun, une boîte englobante ou un contour dans le plan image ; deux détecteurs interchangeables sont étudiés, un seuillage couleur (section~\ref{sec:detection-v1}) et un détecteur appris piloté par le langage (section~\ref{sec:detection-v2}). Puis la \emph{reconstruction 3D} : les détections d'un même objet dans les deux caméras fixes sont appariées, puis \emph{triangulées} pour retrouver sa position dans le repère de base (section~\ref{sec:reconstruction}). Enfin le \emph{filtrage} : les positions hors de l'espace de travail, ou tombant dans une zone d'exclusion connue, sont écartées avant de constituer la scène.

\dansLeProjet{Dans ce travail, la reconstruction 3D repose sur la paire stéréo fixe \texttt{cam\_0}--\texttt{cam\_1} ; la caméra de poignet \texttt{cam\_2} sert surtout au raffinement de la prise (chapitre~\ref{ch:controle}). Le filtrage s'appuie sur une description de la scène qui évite, par exemple, de confondre le filament orange du robot lui-même avec le cube orange visé. Un repli monoculaire par \emph{PnP} est prévu lorsque la stéréo échoue, mais il reste peu opérant sur les formes étudiées --- ce point est discuté à la section~\ref{sec:reconstruction}.}

Au-delà de la position, la perception estime aussi l'\emph{encombrement} de l'objet --- une boîte englobante tridimensionnelle approchée --- ainsi que des indices d'orientation. Ces grandeurs alimentent la planification de saisie, qui en déduit l'ouverture et l'alignement des mâchoires (chapitre~\ref{ch:planification}).

La suite du chapitre détaille ces éléments dans l'ordre où ils se composent : les modèles de caméra et la calibration qui fondent toute la géométrie (section~\ref{sec:calibration-cam}), les deux détecteurs 2D (sections~\ref{sec:detection-v1} et~\ref{sec:detection-v2}), la reconstruction 3D par triangulation et sa sensibilité au bruit (section~\ref{sec:reconstruction}), la synchronisation des caméras (section~\ref{sec:sync}), et enfin la validation expérimentale de la perception (section~\ref{sec:perception-valid}).

\section{Calibration et modèles caméra}
\label{sec:calibration-cam}

Toute la géométrie de la perception repose sur deux ingrédients : un \emph{modèle} qui relie un point de l'espace au pixel où la caméra le voit, et une \emph{calibration} qui en mesure les paramètres pour chaque caméra. Cette section établit d'abord le modèle sténopé et sa distorsion --- les paramètres \emph{intrinsèques}, internes à l'optique ---, puis la calibration œil-main qui place la caméra dans le repère du robot --- les paramètres \emph{extrinsèques} ---, et enfin la chaîne de repères qui en résulte.

\paragraph{Le modèle sténopé.} Une caméra est modélisée comme un \emph{sténopé} (\emph{pinhole}) : tout rayon lumineux passe par un unique centre optique avant de frapper le plan image. Concrètement, la lumière émise par un point de la scène traverse ce centre et vient frapper le \emph{capteur} --- une grille de pixels --- en un seul endroit ; c'est la position de ce point d'impact, en pixels, que le modèle cherche à calculer. Et comme tous les rayons convergent vers le même centre, un objet éloigné se projette sur une plus petite portion du capteur qu'un objet proche : c'est l'effet de \emph{perspective}, que va traduire la division par la profondeur. Un point $P_{\text{cam}} = (X, Y, Z)$ exprimé dans le repère de la caméra --- l'axe $Z$ pointant vers la scène --- se projette alors en deux temps. D'abord une \emph{division perspective}, qui traduit qu'un objet paraît d'autant plus petit qu'il est loin :
\[
  x = \frac{X}{Z}, \qquad y = \frac{Y}{Z},
\]
où $(x, y)$ sont les coordonnées sur le plan image \emph{normalisé}. Ensuite une mise à l'échelle par les distances focales, puis une translation par le point principal pour exprimer le résultat en pixels :
\[
  u = f_x\,x + c_x, \qquad v = f_y\,y + c_y ,
\]
où $(u,v)$ est le pixel obtenu, $f_x$ et $f_y$ les \emph{distances focales} (en pixels) et $(c_x, c_y)$ le \emph{point principal}.
Les deux étapes se regroupent sous forme matricielle, en coordonnées homogènes, à l'aide de la matrice \emph{intrinsèque} $K$ :
\[
  s \begin{bmatrix} u \\ v \\ 1 \end{bmatrix}
  = K \begin{bmatrix} X \\ Y \\ Z \end{bmatrix},
  \qquad
  K = \begin{bmatrix} f_x & 0 & c_x \\ 0 & f_y & c_y \\ 0 & 0 & 1 \end{bmatrix},
  \qquad s = Z .
\]
Le facteur d'échelle $s = Z$ rappelle que la projection \emph{perd} la profondeur : une seule image ne permet pas de retrouver $Z$ --- c'est précisément l'objet de la reconstruction 3D (section~\ref{sec:reconstruction}). Les paramètres de $K$ se lisent physiquement : les distances focales prennent deux valeurs $f_x$ et $f_y$ parce qu'un pixel n'est pas nécessairement carré. Le point principal $(c_x, c_y)$ est l'intersection de l'axe optique avec le plan image : comme les pixels se comptent depuis un \emph{coin} de l'image, ses coordonnées avoisinent la moitié des dimensions --- d'où des valeurs proches de $(960, 540)$ pour une image $1920\times1080$ --- sans pour autant coïncider avec ce centre géométrique, puisqu'il s'agit d'une grandeur \emph{mesurée} par la calibration et non du centre supposé de l'image. Enfin, ces paramètres sont \emph{internes} à la caméra : ils ne dépendent pas de l'endroit où elle se trouve et valent pour toutes ses images.

\dansLeProjet{Dans ce travail, chaque caméra est calibrée à la résolution de $1920\times1080$ pixels. Pour \texttt{cam\_0}, la calibration (décrite plus bas) donne des distances focales $f_x \approx 1226$ et $f_y \approx 1225$ pixels --- des pixels donc quasi carrés --- et un point principal $(c_x, c_y) \approx (979,\,610)$, voisin du centre géométrique de l'image $(960,\,540)$.}

\begin{figure}[ht]
  \centering
  \begin{tikzpicture}[>={Stealth[round]}, font=\small]
    \coordinate (O)     at (0,0);        % centre optique
    \coordinate (Ofoot) at (4.5,0);      % pied de l'objet (sur l'axe)
    \coordinate (Otip)  at (4.5,1.7);    % sommet de l'objet
    \coordinate (c)     at (-2.2,0);     % point principal (sur le capteur)
    \coordinate (Itip)  at (-2.2,-0.831);% image du sommet (inversee) : -1.7*2.2/4.5
    % plan image (capteur), DERRIERE le centre optique
    \fill[blue!8] (-2.45,-1.5) rectangle (-1.95,1.5);
    \draw[thick,blue!55!black] (-2.2,-1.5) -- (-2.2,1.5);
    \node[blue!55!black, above, align=center] at (-2.2,1.5) {plan image\\(capteur)};
    % axe optique, Z vers la scene (droite)
    \draw[->,dashed] (-3.3,0) -- (5.7,0) node[below right] {axe optique $Z$ (vers la scène)};
    % rayons : sommet objet -> O -> image (sens de la lumiere)
    \draw[thick,-{Stealth[round]}] (Otip) -- (O);
    \draw[thick,-{Stealth[round]}] (O) -- (Itip);
    % objet (fleche verticale, a droite)
    \draw[very thick,-{Stealth[round]}] (Ofoot) -- (Otip);
    \node[right=3pt] at (Otip) {objet $P=(X,Y,Z)$};
    % image inversee (fleche vers le bas, sur le capteur)
    \draw[very thick,-{Stealth[round]},red!75!black] (c) -- (Itip);
    \node[red!75!black, left=2pt, align=right] at (Itip) {image\\inversée\\$p=(u,v)$};
    % centre optique
    \fill (O) circle (1.9pt);
    \node[above right=0pt and 1pt] at (O) {$O$ \footnotesize(centre optique)};
    % distance focale f (O -> capteur)
    \draw[<->] (0,-1.95) -- (-2.2,-1.95) node[midway,fill=white,inner sep=1.5pt] {$f$};
    % profondeur Z (O -> objet)
    \draw[<->] (0,2.5) -- (4.5,2.5) node[midway,fill=white,inner sep=1.5pt] {$Z$ (profondeur)};
  \end{tikzpicture}
  \caption{Le modèle de caméra sténopé : la lumière d'un point $P$ de la scène traverse le centre optique $O$ et forme une image \emph{inversée} sur le capteur, placé derrière. La division par la profondeur $Z$ produit l'effet de perspective.}
  \label{fig:pinhole}
\end{figure}

\paragraph{La distorsion.} Le modèle sténopé idéalise l'optique : il suppose un trou ponctuel parfait. Une vraie caméra emploie en réalité une \emph{lentille} --- une pièce de verre \emph{courbe} --- pour capter davantage de lumière, et cette courbure ne dévie pas tous les rayons aussi régulièrement qu'un trou idéal : l'écart grandit à mesure que le rayon frappe la lentille loin de son centre, là où il traverse le verre plus obliquement (c'est-à-dire vers les bords de l'image). Visuellement, les lignes droites de la scène apparaissent alors légèrement \emph{bombées}, surtout près des bords --- c'est la \emph{distorsion}. Sa composante principale est \emph{radiale} (les bords « gonflent » : effet barillet ; ou « se creusent » : effet coussinet) ; s'y ajoute une faible distorsion \emph{tangentielle}, due à un léger défaut de parallélisme entre la lentille et le capteur. On la modélise sur les coordonnées normalisées $(x,y)$, avant le passage en pixels. En posant $r^2 = x^2 + y^2$ (le carré de la distance au centre optique), le modèle employé --- celui d'OpenCV --- s'écrit
\[
\begin{aligned}
  x_d &= x\,(1 + k_1 r^2 + k_2 r^4 + k_3 r^6) + 2 p_1 x y + p_2 (r^2 + 2x^2),\\
  y_d &= y\,(1 + k_1 r^2 + k_2 r^4 + k_3 r^6) + p_1 (r^2 + 2y^2) + 2 p_2 x y,
\end{aligned}
\]
où $(k_1, k_2, k_3)$ gouvernent la part radiale et $(p_1, p_2)$ la part tangentielle. Ce sont les coordonnées \emph{distordues} $(x_d, y_d)$ que la caméra forme réellement ; connaître ces coefficients permet de \emph{redresser} l'image avant tout calcul géométrique.

\paragraph{Calibration intrinsèque : la méthode de Zhang.} La matrice $K$ et les coefficients de distorsion ne se mesurent pas directement : on les estime par la méthode de \emph{Zhang}~\cite{zhang_flexible_2000}. Le principe est de photographier une mire plane connue --- un \emph{damier} --- sous de nombreuses orientations : chaque vue fournit des contraintes (via l'homographie entre le plan du damier et l'image), et leur ensemble permet de résoudre pour $K$ et la distorsion, puis d'affiner ces paramètres en minimisant l'\emph{erreur de reprojection} --- l'écart moyen, en pixels, entre les coins du damier détectés dans l'image et leur reprojection par le modèle estimé. Cette erreur sert aussi de mesure de qualité de la calibration.

\dansLeProjet{Dans ce travail, la calibration \emph{intrinsèque} de chaque caméra utilise un damier de $7\times7$ coins internes (carrés de $22{,}2$ mm) présenté sous une quarantaine de poses. Pour \texttt{cam\_0}, le coefficient radial dominant vaut $k_1 \approx -0{,}41$ (une distorsion en barillet), les termes tangentiels étant négligeables (capteur bien aligné), et l'erreur de reprojection finale n'est que de $0{,}19$ pixel --- une précision \emph{sous-pixellique}. La qualité de calibration des trois caméras est reprise à la section~\ref{sec:perception-valid}.}

\begin{figure}[ht]
  \centering
  \photoplaceholder[0.7\linewidth]{Photo à insérer : une vue de la calibration intrinsèque (le damier $7\times7$ présenté à une caméra). Une seule illustration suffit, la procédure étant identique pour les trois caméras.}
  \caption{Calibration intrinsèque par la méthode de Zhang : le damier est présenté à la caméra sous différentes orientations. La même procédure s'applique aux trois caméras.}
  \label{fig:calib-intrinseque}
\end{figure}

\paragraph{Calibration œil-main : où est la caméra par rapport au robot ?} Détecter un objet dans une image ne donne qu'un \emph{pixel} ; pour que le robot agisse dessus, il lui faut savoir \emph{où se trouve la caméra par rapport à lui}. Sans cela, voir l'objet ne sert à rien : on ignorerait dans quelle direction tendre le bras. Déterminer la position et l'orientation exactes de chaque caméra dans le repère de base du robot, c'est l'objet de la calibration \emph{œil-main} (\emph{hand-eye}), dont le principe a été posé au chapitre~\ref{ch:etat} et que l'on démontre ici sur les repères du système.

On ne peut pas relever cette pose à la règle --- l'orientation, surtout, exige une précision qu'aucune mesure manuelle n'atteint. On la \emph{calibre} donc, en montrant à la caméra une mire connue --- un \emph{damier} --- dans de nombreuses poses du bras. Chaque prise fournit alors deux transformations, et deux seulement : la pose du bras $T_{\text{base}\leftarrow\text{pince}}$, donnée par la \emph{cinématique directe} (chapitre~\ref{ch:controle}), et la pose du damier vue par la caméra $T_{\text{cam}\leftarrow\text{damier}}$ --- où $\text{cam}$ désigne la caméra en cours de calibration ---, obtenue en détectant les coins du damier dans l'image puis en résolvant le problème de perspective (\emph{PnP}, fonction \texttt{solvePnP}). Tout l'enjeu est de combiner ces deux mesures pour en extraire la pose, inconnue, de la caméra ; et la manière de les combiner dépend du \emph{montage} (les deux cas sont schématisés en figure~\ref{fig:handeye-loops}).

Reste à expliquer comment se \emph{mesure} cette seconde pose, $T_{\text{cam}\leftarrow\text{damier}}$ : c'est le modèle sténopé du début de section, pris \emph{à l'envers}. Ce modèle envoyait un point connu du repère caméra vers un pixel ; notons $\pi_K$ cette projection complète (division perspective, distorsion, puis passage en pixels par $K$). Ici on fait l'inverse : on connaît les pixels et l'on cherche la pose. Chaque coin du damier a une position 3D connue $M_k$ (la mire est une grille plane régulière), se projette en un pixel $m_k$ que l'on détecte dans l'image, et $K$ a déjà été calibrée ; la seule inconnue est donc la pose qui, glissée dans $\pi_K$, fait retomber les coins projetés sur les pixels observés. On retient celle qui minimise leur écart total --- c'est-à-dire exactement l'\emph{erreur de reprojection} déjà rencontrée pour la calibration de Zhang, mais à $K$ \emph{figé} :
\[
  T_{\text{cam}\leftarrow\text{damier}} \;=\; \arg\min_{T}\; \sum_k \big\lVert\, \pi_K\!\big(T\,M_k\big) - m_k \,\big\rVert^2 .
\]
Avec $6\times9 = 54$ coins pour seulement six inconnues de pose, le problème est très surdéterminé ; c'est le problème classique \emph{Perspective-n-Point} (PnP), que l'on résout par \texttt{solvePnP}. On dispose ainsi, à chaque prise, des deux transformations voulues.

\begin{figure}[t]
  \centering
  \begin{subfigure}{\linewidth}
    \centering
    \begin{tikzpicture}[
        >={Stealth[round]},
        every node/.style={font=\small},
        frame/.style={draw, rounded corners=2pt, minimum width=16mm, minimum height=8mm},
        meas/.style={->, thick},
        cible/.style={->, very thick, dashed, red!75!black},
        aux/.style={->, thick, dashed, gray!55!black},
        elab/.style={font=\small, inner sep=1.5pt, fill=white, pos=0.5}
      ]
      \node[frame] (base)  at (0,2)   {base};
      \node[frame] (pince) at (4.2,2) {pince};
      \node[frame] (cam)   at (4.2,0) {\texttt{cam\_2}};
      \node[frame] (dam)   at (0,0)   {damier};
      \draw[meas]  (base)  -- (pince) node[elab, above]{$T_{\text{base}\leftarrow\text{pince}}$};
      \draw[cible] (pince) -- (cam)   node[elab, right]{$X = T_{\text{pince}\leftarrow\text{cam2}}$};
      \draw[meas]  (cam)   -- (dam)   node[elab, below]{$T_{\text{cam2}\leftarrow\text{damier}}$};
      \draw[aux]   (base)  -- (dam)   node[elab, left] {$T_{\text{base}\leftarrow\text{damier}}$};
    \end{tikzpicture}
    \caption{\emph{Eye-in-hand} (\texttt{cam\_2}) : le damier est fixe sur la table, donc $T_{\text{base}\leftarrow\text{damier}}$ est une constante que l'on élimine ; l'unique inconnue cherchée est $X$.}
    \label{fig:loop-eih}
  \end{subfigure}

  \vspace{5mm}

  \begin{subfigure}{\linewidth}
    \centering
    \begin{tikzpicture}[
        >={Stealth[round]},
        every node/.style={font=\small},
        frame/.style={draw, rounded corners=2pt, minimum width=16mm, minimum height=8mm},
        meas/.style={->, thick},
        cible/.style={->, very thick, dashed, red!75!black},
        aux/.style={->, thick, dashed, gray!55!black},
        elab/.style={font=\small, inner sep=1.5pt, fill=white, pos=0.5}
      ]
      \node[frame] (base)  at (0,2)   {base};
      \node[frame] (pince) at (4.2,2) {pince};
      \node[frame] (dam)   at (4.2,0) {damier};
      \node[frame] (cam)   at (0,0)   {\texttt{cam\_0}};
      \draw[meas]  (base)  -- (pince) node[elab, above]{$T_{\text{base}\leftarrow\text{pince}}$};
      \draw[aux]   (pince) -- (dam)   node[elab, right]{$Z = T_{\text{pince}\leftarrow\text{damier}}$};
      \draw[meas]  (cam)   -- (dam)   node[elab, below]{$T_{\text{cam0}\leftarrow\text{damier}}$};
      \draw[cible] (base)  -- (cam)   node[elab, left] {$Y = T_{\text{base}\leftarrow\text{cam0}}$};
    \end{tikzpicture}
    \caption{\emph{Eye-to-hand} (\texttt{cam\_0}) : le damier est porté par la pince ; il y a donc \emph{deux} constantes inconnues, la cible $Y$ et l'auxiliaire $Z$ que l'on élimine.}
    \label{fig:loop-eth}
  \end{subfigure}

  \caption{Les deux montages vus comme des \emph{boucles de repères}. Chaque flèche de $A$ vers $B$ porte la transformation $T_{A\leftarrow B}$ (la pose de $B$ dans $A$, chapitre~\ref{ch:architecture}). Dans chaque boucle, les \emph{deux} chemins menant au damier sont égaux : cette égalité, écrite pour deux poses du bras, fournit l'équation $A\,X = X\,B$. \textbf{Trait noir plein} : grandeur mesurée à chaque pose --- $T_{\text{base}\leftarrow\text{pince}}$ par cinématique directe, $T_{\text{cam}\leftarrow\text{damier}}$ par PnP. \textbf{Trait tireté} : transformation \emph{constante} inconnue --- en rouge la pose de caméra cherchée, en gris l'auxiliaire éliminée par le calcul.}
  \label{fig:handeye-loops}
\end{figure}

\paragraph{Premier montage : la caméra de poignet (\emph{eye-in-hand}).} La caméra \texttt{cam\_2} est montée sur le bras et se déplace avec lui (montage schématisé en figure~\ref{fig:loop-eih}). On pose donc le damier \emph{fixe} sur la table, et l'on promène le bras pour le lui faire voir sous des angles variés. L'inconnue cherchée est la pose, \emph{constante}, de la caméra par rapport à la pince, $X = T_{\text{pince}\leftarrow\text{cam2}}$. Pour l'isoler, on écrit qu'à chaque pose $i$ le damier occupe la \emph{même} place dans le repère de base --- il ne bouge pas ---, en parcourant la chaîne $\text{base}\to\text{pince}\to\text{cam2}\to\text{damier}$ :
\[
  T_{\text{base}\leftarrow\text{damier}} \;=\; T_{\text{base}\leftarrow\text{pince}}(i)\; X\; T_{\text{cam2}\leftarrow\text{damier}}(i) ,
\]
le membre de gauche étant une constante (la place du damier sur la table). En égalant cette relation pour deux poses $i$ et $j$, cette constante s'élimine ; après réarrangement, il reste
\[
  \underbrace{T_{\text{base}\leftarrow\text{pince}}(j)^{-1}\,T_{\text{base}\leftarrow\text{pince}}(i)}_{A\;:\;\text{déplacement de la pince}}\; X
  \;=\; X\;
  \underbrace{T_{\text{cam2}\leftarrow\text{damier}}(j)\,T_{\text{cam2}\leftarrow\text{damier}}(i)^{-1}}_{B\;:\;\text{déplacement de la caméra}} ,
\]
c'est-à-dire l'équation de calibration œil-main
\[
  A\,X = X\,B .
\]
Son sens géométrique se lit directement : entre les poses $i$ et $j$, la pince effectue un mouvement $A$ et la caméra un mouvement $B$ ; comme la caméra est rigidement vissée à la pince, c'est le \emph{même} déplacement décrit dans deux repères, et la transformation constante $X$ qui passe de l'un à l'autre est précisément ce qui les relie.

Une seule paire de poses ne suffit pas à fixer $X$ --- une transformation rigide compte six degrés de liberté ---, et sa rotation n'est contrainte que si les poses tournent autour d'axes \emph{différents}. On accumule donc des dizaines de poses bien variées, ce qui \emph{surdétermine} le système $A\,X = X\,B$. La méthode de \emph{Horaud}~\cite{horaud_hand-eye_1995}, exposée par la fonction \texttt{calibrateHandEye} d'OpenCV, en estime la meilleure solution au sens des moindres carrés en optimisant \emph{conjointement} la rotation et la translation de $X$, ce qui la rend robuste au bruit de mesure (chapitre~\ref{ch:etat}). Concrètement, on \emph{empile} l'équation $A\,X = X\,B$ pour toutes les paires de poses et l'on résout ce système surdéterminé : sa résolution numérique relève de la méthode citée et n'est pas reprise ici, tandis que la \emph{qualité} du $X$ obtenu --- l'écart résiduel sur l'ensemble des poses --- est, elle, mesurée à la section~\ref{sec:perception-valid}. Pour \texttt{cam\_2}, la méthode renvoie la matrice constante $T_{\text{pince}\leftarrow\text{cam2}}$.

La caméra de poignet bougeant avec le bras, sa pose dans la base \emph{change} à chaque instant ; on la reconstruit à la volée en composant la pose courante de la pince (cinématique directe) avec la transformation $X$ que l'on vient de calibrer :
\[
  T_{\text{base}\leftarrow\text{cam2}}(t) \;=\; T_{\text{base}\leftarrow\text{pince}}(t)\; \underbrace{T_{\text{pince}\leftarrow\text{cam2}}}_{=\,X} .
\]

\paragraph{Second montage : les caméras fixes (\emph{eye-to-hand}).} Les deux caméras de la paire stéréo, \texttt{cam\_0} et \texttt{cam\_1}, sont au contraire vissées à la structure : elles ne bougent jamais (montage schématisé en figure~\ref{fig:loop-eth}). Le damier est cette fois collé \emph{sur la pince}, et c'est le bras qui le promène devant la caméra immobile. L'inconnue recherchée est alors \emph{directement} la pose de la caméra dans la base, $Y = T_{\text{base}\leftarrow\text{cam0}}$. Les rôles sont simplement échangés par rapport au premier montage : c'est désormais le damier (porté par la pince) qui bouge, tandis que la caméra reste fixe. Deux trajets mènent alors au damier --- par la caméra fixe, ou par le bras --- et doivent coïncider :
\[
  \underbrace{Y}_{T_{\text{base}\leftarrow\text{cam0}}}\; T_{\text{cam0}\leftarrow\text{damier}}(i)
  \;=\; T_{\text{base}\leftarrow\text{pince}}(i)\; \underbrace{Z}_{T_{\text{pince}\leftarrow\text{damier}}} ,
\]
où $Z$, la pose du damier sur la pince, est une seconde constante inconnue --- auxiliaire, dont on ne veut pas. Pour l'écarter, on l'isole,
\[
  Z = T_{\text{base}\leftarrow\text{pince}}(i)^{-1}\; Y\; T_{\text{cam0}\leftarrow\text{damier}}(i) ,
\]
puis, $Z$ ne dépendant pas de la pose, on égale ses deux expressions en $i$ et $j$ ; après réarrangement, l'auxiliaire a disparu et il reste
\[
  \underbrace{T_{\text{base}\leftarrow\text{pince}}(j)\,T_{\text{base}\leftarrow\text{pince}}(i)^{-1}}_{A'}\; Y
  \;=\; Y\;
  \underbrace{T_{\text{cam0}\leftarrow\text{damier}}(j)\,T_{\text{cam0}\leftarrow\text{damier}}(i)^{-1}}_{B'} .
\]
C'est, terme pour terme, la \emph{même} équation $A\,X = X\,B$ que pour la caméra de poignet --- à une seule différence : la pose du bras y entre \emph{inversée} (comparer $A'$ à $A$ : $T_{\text{base}\leftarrow\text{pince}}^{-1} = T_{\text{pince}\leftarrow\text{base}}$). Il suffit donc de fournir au même solveur de Horaud les poses du bras \emph{inversées} ; il renvoie alors directement la pose cherchée, $Y = T_{\text{base}\leftarrow\text{cam0}}$. C'est, à la ligne près, ce qu'opère le code (\texttt{src/calibration/handeye.py}) : inverser la pose du bras avant l'appel à \texttt{calibrateHandEye}, puis relire sa sortie comme une pose caméra-vers-base. Le résultat est une matrice \emph{constante}, réutilisée telle quelle à chaque image. \texttt{cam\_1} possède, elle, ses propres paramètres \emph{intrinsèques}, mais sa \emph{pose} n'est pas estimée par un hand-eye indépendant : on résout d'abord \texttt{cam\_0} seule, puis l'on \emph{déduit} celle de \texttt{cam\_1} en la composant avec la transformation stéréo $T_{\text{cam0}\leftarrow\text{cam1}}$ (mesurée par calibration stéréo du couple). Ce choix garantit la \emph{cohérence} de la paire par construction --- essentielle à la triangulation (section~\ref{sec:perception-valid}).

\dansLeProjet{Dans ce travail, le damier de calibration extrinsèque est volontairement \emph{asymétrique} --- $6\times9$ coins internes (carrés de $22$ mm), soit un nombre \emph{pair} de coins dans un sens et \emph{impair} dans l'autre ---, ce qui lève l'ambiguïté d'orientation à $180^\circ$ d'un damier symétrique et fiabilise l'estimation PnP. Chaque caméra est calibrée par la méthode de Horaud sur plusieurs dizaines de poses du bras (jusqu'à $71$ capturées), dont une partie seulement est retenue après rejet automatique des prises mal détectées ; le décompte par caméra et la cohérence des poses obtenues sont reportés à la section~\ref{sec:perception-valid}.}

\begin{figure}[ht]
  \centering
  \photoplaceholder[0.85\linewidth]{Photos à insérer : les deux montages de calibration œil-main. \emph{Eye-to-hand} (\texttt{cam\_0}/\texttt{cam\_1}) : le damier est fixé sur la pince et le bras bouge devant la caméra fixe. \emph{Eye-in-hand} (\texttt{cam\_2}) : le damier est posé sur la table et la caméra, portée par le bras, l'observe sous divers angles.}
  \caption{Calibration œil-main : montage \emph{eye-to-hand} (damier sur la pince) et \emph{eye-in-hand} (damier sur la table). Contrairement à l'intrinsèque, le montage diffère selon le rôle de la caméra.}
  \label{fig:calib-extrinseque}
\end{figure}

\paragraph{Mettre bout à bout : la matrice de projection.} On dispose maintenant des deux pièces : la matrice intrinsèque $K$, qui projette un point du repère \emph{caméra} sur un pixel (modèle sténopé, début de section), et la pose extrinsèque $T_{\text{base}\leftarrow\text{cam}}$, qui situe la caméra dans la base (calibration œil-main). En les composant, on obtient un \emph{unique} opérateur qui envoie directement un point du repère \emph{base} sur un pixel --- c'est ce dont la reconstruction 3D aura besoin.

Suivons un point $\mathbf{X} = (X, Y, Z, 1)^{\top}$ du repère base, en coordonnées homogènes. On l'amène d'abord dans le repère caméra par l'inverse de la pose, $T_{\text{cam}\leftarrow\text{base}} = (T_{\text{base}\leftarrow\text{cam}})^{-1}$, puis on le projette par $K$. Or cette pose est une matrice $4\times4$ de la forme $\left[\begin{smallmatrix} R & t\\ 0 & 1 \end{smallmatrix}\right]$ ; en n'en gardant que les trois premières lignes --- le bloc $3\times4$ noté $[\,R \mid t\,]$ ---, les deux étapes se condensent en une seule \emph{matrice de projection} :
\[
  s \begin{bmatrix} u \\ v \\ 1 \end{bmatrix}
  \;=\; \underbrace{K\,[\,R \mid t\,]}_{P}\; \mathbf{X},
  \qquad [\,R \mid t\,] = T_{\text{cam}\leftarrow\text{base}}\ \text{(sans sa dernière ligne)} .
\]
La matrice $P$, de taille $3\times4$, résume ainsi \emph{toute} la géométrie d'une caméra --- intrinsèque \emph{et} pose --- en un seul opérateur. Elle n'en décrit toutefois que la partie \emph{linéaire} : la distorsion, non linéaire (début de section), est retirée des points \emph{avant} d'appliquer $P$ (section~\ref{sec:reconstruction}). Pour les deux caméras fixes, $P$ est \emph{constante} ; pour la caméra de poignet, elle se recalcule à chaque instant à partir de la pose courante de la pince (calibration œil-main ci-dessus).

C'est précisément cet opérateur que la reconstruction 3D va \emph{inverser} : avec deux caméras --- donc deux matrices $P$ --- on remonte d'un couple de pixels à la position 3D d'un point, par \emph{triangulation} (section~\ref{sec:reconstruction}).

\section{Détection 2D --- version V1 (HSV)}
\label{sec:detection-v1}

La détection 2D localise, dans chaque image, les objets connus et en renvoie pour chacun une \texttt{Detection2D} --- un label, un centre, une boîte englobante, un contour et un score (chapitre~\ref{ch:architecture}). Deux détecteurs interchangeables sont étudiés derrière la même interface ; cette section présente le premier, \emph{déterministe} et fondé sur la couleur~\cite{forsyth_computer_2012}.

Son rôle n'est pas d'être le plus puissant, mais d'\emph{isoler la contribution de la géométrie}. Parce qu'il est reproductible et sans apprentissage, toute erreur de position 3D mesurée en aval ne peut venir que de la calibration ou de la triangulation --- jamais d'un réseau dont le comportement varierait d'une exécution à l'autre. C'est donc le détecteur de référence sur lequel reposent la mise au point et la validation (section~\ref{sec:perception-valid}).

\paragraph{Pourquoi l'espace HSV.} En RGB, la couleur d'un objet se déplace sur les trois canaux à la fois dès que l'éclairage change, ce qui rend tout seuillage fragile. L'espace \emph{HSV} (\emph{teinte}, \emph{saturation}, \emph{valeur}) sépare au contraire la \emph{chrominance} --- la teinte $H$ et la saturation $S$ --- de la \emph{luminance} --- la valeur $V$. Un objet rouge garde à peu près la même teinte à l'ombre ou en pleine lumière ; seule sa valeur change. Seuiller la couleur en HSV est donc bien plus stable face à l'éclairage qu'en RGB. Un troisième espace, \emph{HSL} (\emph{teinte}, \emph{saturation}, \emph{luminosité}), partage ce mérite --- lui aussi isole la teinte ---, mais on lui préfère HSV pour deux raisons. D'une part, la \emph{valeur} $V = \max(R,G,B)$ épouse directement l'exposition de la caméra, là où la \emph{luminosité} $L = \tfrac{1}{2}(\max + \min)$ de HSL place une couleur \emph{pure et vive} à mi-échelle ($L \approx 0{,}5$) et mêle ainsi « vivacité » et « clarté », moins commodes à seuiller séparément. D'autre part, la saturation HSV sépare nettement une couleur vive d'un blanc --- un pixel presque blanc y est peu saturé ---, tandis que la saturation HSL attribue au contraire une forte valeur aux teintes très \emph{claires}, ce qui brouille le seuillage. Chaque image est donc d'abord convertie en HSV.

\paragraph{De l'image au contour.} Chaque classe d'objet est décrite par une \texttt{ObjectSpec} : un label (\texttt{red\_cube}, \texttt{green\_cylinder}\dots), une plage HSV et des bornes d'aire. Pour une classe donnée, le détecteur procède en quatre temps. Il construit d'abord un \emph{masque} binaire, qui retient les pixels dont le triplet $(H,S,V)$ tombe dans la plage. Une étape de \emph{morphologie} --- une ouverture, qui efface les points isolés, suivie d'une fermeture, qui bouche les trous --- nettoie ce masque. On en extrait alors les \emph{contours} des régions connexes, dont on ne garde que celles d'aire comprise entre un minimum (rejet du bruit) et un maximum (rejet des aplats parasites, tel un fond mal éclairé). De chaque contour retenu, enfin, on tire le \emph{centre} --- le centroïde, calculé par les moments de l'image ---, la boîte englobante et un score (l'aire, normalisée) ; on ne conserve que les $k$ plus grands par classe, par défaut un seul, en supposant un exemplaire visible (figure~\ref{fig:detection-hsv}).

\begin{figure}[t]
  \centering
  \photoplaceholder[0.9\linewidth]{Photo à insérer : les étapes du détecteur HSV sur une prise réelle --- (a) image d'origine, (b) masque binaire après seuillage HSV et nettoyage morphologique, (c) contour retenu et son centroïde --- sur un objet coloré posé sur la table.}
  \caption{Détection HSV, du pixel au contour : seuillage couleur, nettoyage morphologique, puis extraction du contour et de son centre.}
  \label{fig:detection-hsv}
\end{figure}

\paragraph{Le cas des couleurs achromatiques.} Un détail compte plus qu'il n'y paraît : la teinte n'a \emph{aucun sens} pour le noir, le blanc ou le gris. Seuiller un objet noir comme une teinte capterait indistinctement toutes les couleurs sombres. Ces cas sont donc traités à part --- le noir par sa seule valeur ($V$ faible, $H$ et $S$ ignorés), le blanc par une faible saturation jointe à une valeur élevée ---, ce qui évite des confusions massives entre objets sombres, ou entre objets clairs. De même, le rouge chevauchant la discontinuité de la teinte à $H = 0/179$ (convention OpenCV, où la teinte est codée de $0$ à $179$), il est décrit par \emph{deux} plages réunies.

\dansLeProjet{Dans ce travail, les plages couleur des primitives sont réglées sous l'éclairage du poste et stockées dans \texttt{configs/perception/hsv\_specs.json}. Le détecteur autorise des plages \emph{par caméra} : la caméra de poignet (\texttt{cam\_2}), en vue rapprochée et souvent plus délavée, ne répond pas tout à fait comme la paire stéréo fixe --- un même objet y paraît moins saturé ---, et distinguer les plages par caméra reste cohérent avec une calibration déjà propre à chacune.}

\paragraph{Forces et limites.} Ce détecteur est \emph{rapide}, \emph{reproductible} et \emph{sans dépendance d'apprentissage} --- exactement ce qu'on attend d'un socle. Mais il ne « voit » que de la couleur : il signale « il y a des pixels oranges ici » sans savoir \emph{ce que} c'est. Il ne distingue donc pas le cube orange visé des parties oranges du robot lui-même (d'où le filtrage de scène, section~\ref{sec:perception-pipeline}) ; il exige une calibration des plages sous l'éclairage d'usage ; et il reste aveugle à la \emph{forme} et au \emph{contexte}. Ces limites motivent le second détecteur --- \emph{ouvert} et piloté par le langage --- présenté à la section suivante.

\section{Détection 2D --- version V2 (\textit{open-vocabulary})}
\label{sec:detection-v2}

Le détecteur couleur a un angle mort : il ne raisonne que sur des pixels, jamais sur le \emph{sens} de ce qu'il voit. Le second détecteur lève cette limite en s'appuyant sur un modèle \emph{vision-langage} pré-entraîné, capable de localiser un objet décrit en \emph{langage naturel} --- sans liste de classes figée ni couleur réglée d'avance. C'est la détection à \emph{vocabulaire ouvert} (\emph{open-vocabulary}).

\paragraph{Le principe.} On fournit au modèle, en plus de l'image, une \emph{liste de descriptions textuelles} --- « \emph{a small bright orange plastic cube} », « \emph{a purple plastic cylinder} »\dots --- et il renvoie, pour chacune, les boîtes englobantes correspondantes assorties d'un score. Entraîné sur des millions de paires image--légende, il a appris à relier les mots à l'apparence : il peut alors séparer le cube orange des \emph{parties oranges du robot} lui-même par la \emph{forme} et le \emph{contexte}, et non par la seule couleur comme le faisait HSV. Changer les objets cibles ne demande alors que de réécrire les descriptions --- ni recalibration, ni réentraînement.

\paragraph{Le modèle retenu.} Parmi les détecteurs à vocabulaire ouvert recensés au chapitre~\ref{ch:etat}, on retient \emph{OWL-ViTv2}~\cite{minderer_owlv2_2023} (\texttt{google/owlv2-base-patch16-ensemble}), servi par la bibliothèque \emph{Transformers} de Hugging Face. Ce choix est cohérent avec l'écosystème du projet : c'est le même \emph{écosystème Hugging Face} (outillage et \emph{Hub} de modèles) que celui des politiques apprises de LeRobot, employées par la seconde approche de ce travail --- aucune fragmentation des dépendances.

\paragraph{De la sortie du modèle à une \texttt{Detection2D}.} Chaque boîte renvoyée au-dessus d'un seuil de confiance devient une \texttt{Detection2D} : son centre est le milieu de la boîte, et son « contour » se réduit aux quatre coins de celle-ci. Ce pseudo-contour --- faute de segmentation fine --- suffit à une prise par le dessus, mais reste trop grossier pour fixer précisément l'orientation du poignet (chapitre~\ref{ch:planification}). La sortie respecte exactement la même interface que le détecteur HSV : tout l'aval du pipeline est identique, qu'on emploie l'un ou l'autre (figure~\ref{fig:detection-owl}).

\begin{figure}[t]
  \centering
  \photoplaceholder[0.9\linewidth]{Photo à insérer : une détection \emph{open-vocabulary} (OWL-ViTv2) sur une scène réelle, avec les boîtes et leurs labels textuels --- par exemple « \emph{orange cube} » distingué du bras du robot, tous deux oranges, ce que le seuillage HSV ne saurait faire.}
  \caption{Détection à vocabulaire ouvert : le modèle localise chaque objet à partir de sa \emph{description} textuelle et sépare des objets de même couleur par la forme et le contexte.}
  \label{fig:detection-owl}
\end{figure}

\dansLeProjet{Dans ce travail, OWL-ViTv2 (variante \emph{base}, de l'ordre de la centaine de millions de paramètres) tourne sur le GPU intégré du MacBook M4 (\emph{backend} MPS), mais nettement plus lentement que le seuillage HSV --- de l'ordre de quelques secondes par image ---, ce qui le réserve à une détection \emph{ponctuelle} plutôt qu'en flux continu. Les descriptions sont \emph{enrichies} (couleur, forme, matière) plutôt que minimales, et le seuil de confiance a été porté à $0{,}25$ (dans \texttt{configs/perception/hf\_specs.json}) : à $0{,}15$, le modèle étiquetait parfois les parties oranges du robot comme le cube orange (score $\sim 0{,}2$), d'où des reconstructions 3D erronées ; le seuil plus strict écarte ces confusions.}

\paragraph{Forces et limites.} Ce détecteur est \emph{ouvert}, \emph{sémantique} et \emph{sans calibration couleur} --- mais c'est un réseau lourd ($\sim$\,une centaine de millions de paramètres, quelques \emph{secondes} par image --- bien plus lent que le seuillage HSV), moins transparent qu'un seuillage et porteur des incertitudes propres à un détecteur appris (détections manquées ou parasites selon la formulation et l'éclairage). D'où la complémentarité des deux versions : HSV comme socle reproductible pour \emph{isoler} la géométrie, OWL-ViTv2 pour la \emph{robustesse sémantique}. Leur comparaison chiffrée sur la tâche complète figure au chapitre~\ref{ch:evaluation}.

\section{Reconstruction 3D}
\label{sec:reconstruction}

Une détection ne donne qu'un \emph{pixel}, et la projection a perdu la profondeur (le facteur $s = Z$, section~\ref{sec:calibration-cam}). Pour situer l'objet \emph{en 3D} dans le repère de base, il faut deux points de vue : c'est l'objet de la \emph{reconstruction}, qui \emph{inverse} la matrice de projection $P$ posée plus haut.

\paragraph{Triangulation.} À travers sa matrice $P$, un pixel ne désigne pas un point mais une \emph{droite} : tous les points 3D qui se projettent sur lui. Un même objet vu par les deux caméras fixes (\texttt{cam\_0}, \texttt{cam\_1}) définit donc deux droites, qui se coupent --- en théorie --- en sa position. Le bruit de détection fait qu'elles ne se rencontrent pas tout à fait ; on cherche alors le point qui en minimise l'écart. C'est la \emph{triangulation} par \emph{transformation linéaire directe} (DLT)~\cite{hartley_multiple_2018}, fournie par \texttt{cv2.triangulatePoints} à partir de $P_0$ et $P_1$. Comme $P$ ne décrit que la partie linéaire de la projection, les pixels sont d'abord \emph{dé-distordus} (section~\ref{sec:calibration-cam}) avant la triangulation.

La mise en équations est directe. Si $\mathbf{p}_1^{\top}, \mathbf{p}_2^{\top}, \mathbf{p}_3^{\top}$ sont les trois lignes de $P$ et $(u, v)$ le pixel observé, le fait que $\mathbf{X}$ se projette en $(u,v)$ s'écrit, après élimination du facteur d'échelle, en deux équations \emph{linéaires} en $\mathbf{X}$ :
\[
  u\,\mathbf{p}_3^{\top}\mathbf{X} - \mathbf{p}_1^{\top}\mathbf{X} = 0 , \qquad
  v\,\mathbf{p}_3^{\top}\mathbf{X} - \mathbf{p}_2^{\top}\mathbf{X} = 0 .
\]
Ces deux relations valent pour \emph{une} caméra ; la seconde en fournit deux autres, de même forme (avec sa propre matrice $P$ et son pixel). Les deux caméras empilent ainsi \emph{quatre} équations pour les trois libertés de $\mathbf{X}$ : le système homogène surdéterminé $A\,\mathbf{X} = \mathbf{0}$ se résout par décomposition en valeurs singulières, le point cherché étant le vecteur singulier de plus petite valeur. La figure~\ref{fig:triangulation} illustre cette géométrie --- et pourquoi la profondeur y est la direction la moins bien contrainte.

\begin{figure}[t]
  \centering
  \begin{tikzpicture}[>={Stealth[round]}, font=\small]
    \coordinate (C0) at (0,0);
    \coordinate (C1) at (2.2,0);
    \coordinate (X)  at (1.1,4.0);
    % centres optiques
    \fill (C0) circle (2pt); \node[left=2pt] at (C0) {\texttt{cam\_0}};
    \fill (C1) circle (2pt); \node[right=2pt] at (C1) {\texttt{cam\_1}};
    % ligne de base
    \draw[|<->|] (0,-0.55) -- node[fill=white, inner sep=1.5pt]{base $b$} (2.2,-0.55);
    % lignes de visee
    \draw[thick,->] (C0) -- (X);
    \draw[thick,->] (C1) -- (X);
    % region d'incertitude, allongee selon la profondeur (verticale ici)
    \fill[red!22, draw=red!65] (X) ellipse [x radius=2mm, y radius=8mm];
    \node[right=3mm] at (X) {$\mathbf{X}$};
    % annotation profondeur
    \node[red!65!black, align=left, font=\footnotesize] at (3.4,2.2) {incertitude\\ en profondeur};
    \draw[red!65!black,->,shorten >=1pt] (3.35,2.5) to[out=120,in=-20] (1.35,3.7);
  \end{tikzpicture}
  \caption{Triangulation stéréo. Chaque pixel définit, via sa matrice $P$, une \emph{ligne de visée} ; leur intersection donne le point $\mathbf{X}$. Comme la base $b$ est petite devant la profondeur, les deux droites se coupent sous un angle \emph{rasant} : la région compatible avec les deux mesures est \emph{allongée} le long de l'axe caméra, si bien que la profondeur est estimée moins précisément que les directions latérales.}
  \label{fig:triangulation}
\end{figure}

\paragraph{Apparier avant de trianguler.} Encore faut-il savoir \emph{quelles} détections des deux caméras correspondent au même objet. L'appariement se fait par \emph{label} (même classe), puis par score ; chaque paire candidate est \emph{validée} en re-projetant le point 3D reconstruit dans les deux images et en vérifiant que l'écart aux pixels détectés reste sous un seuil (quelques dizaines de pixels). Une paire qui échoue --- objet hors de l'espace de travail, ou écart trop grand --- est écartée au profit de la suivante.

\paragraph{La profondeur, direction fragile.} La triangulation n'est pas également précise dans toutes les directions. Pour deux caméras de \emph{base} $b$ (l'écart entre leurs centres) et de focale $f$, une erreur de $\Delta d$ pixels sur la position relative d'un point se répercute sur la profondeur estimée $z$ selon
\[
  \Delta z \;\approx\; \frac{z^{2}}{b\,f}\,\Delta d ,
\]
qui croît avec le \emph{carré} de la distance et décroît avec la base. Autrement dit : plus l'objet est loin, ou plus les deux caméras sont rapprochées, plus la profondeur est incertaine --- c'est la direction la plus fragile de la reconstruction, là où les deux droites deviennent presque parallèles.

\dansLeProjet{Dans ce travail, la base stéréo \texttt{cam\_0}--\texttt{cam\_1} vaut environ $103$ mm ; aux distances de travail ($30$ à $50$ cm), la profondeur est l'axe le plus bruité. Or la paire fixe \emph{surplombe} la table --- sa ligne de visée est inclinée d'une trentaine de degrés vers le bas ---, si bien que cet axe faible se projette en bonne part sur la \emph{verticale} : la \emph{hauteur} de l'objet est donc l'une des grandeurs les moins précises, ce qui contribue à l'imprécision en hauteur constatée lors des prises (discutée au chapitre~\ref{ch:evaluation}). Deux garde-fous existent : la \emph{validation par reprojection} ci-dessus, et le \emph{raffinement} de la prise par la caméra de poignet juste avant la descente (chapitre~\ref{ch:controle}). La covariance de position, prévue dans le modèle de données, n'est pas propagée explicitement : la fiabilité d'un point est jugée \emph{a posteriori} par son erreur de reprojection.}

\paragraph{Repli monoculaire.} Lorsqu'un objet n'est vu que par une seule caméra, la stéréo est impossible ; un repli \emph{monoculaire} estime alors sa pose par PnP~\cite{lepetit_epnp_2009} --- la même méthode éprouvée que celle de la calibration (section~\ref{sec:calibration-cam}) ---, en exploitant les dimensions \emph{connues} de l'objet. La difficulté ne tient pas à PnP, mais à la géométrie \emph{monoculaire} elle-même : la pose d'un objet quasi \emph{plan} vu de face --- une face de cube, un prisme posé --- est intrinsèquement \emph{ambiguë}, deux poses symétriques expliquant presque aussi bien l'image, quel que soit le solveur. La stéréo lève cette ambiguïté en croisant deux points de vue ; le repli monoculaire ne joue donc qu'un rôle de filet de sécurité sur les formes étudiées, la reconstruction reposant avant tout sur la paire stéréo.

Le résultat de l'étape est une \texttt{ObjectInstance} : la position 3D de l'objet dans le repère de base, son encombrement (boîte englobante 3D) et, le cas échéant, une orientation --- les grandeurs que la planification de saisie va consommer (chapitre~\ref{ch:planification}).

\section{Synchronisation multi-caméras}
\label{sec:sync}

La triangulation suppose que les deux caméras voient la scène \emph{au même instant} : si l'objet --- ou le bras, pour la caméra de poignet --- a bougé entre les deux prises, les droites ne se croisent plus au bon endroit et la position 3D est faussée. La capture doit donc être \emph{synchronisée}.

\paragraph{Capture en deux temps.} Le pilote \texttt{MultiCamera} sépare la \emph{saisie} de l'image de son \emph{décodage} : il déclenche d'abord un \texttt{grab} sur les caméras en \emph{rafale serrée}, puis ne décode (\texttt{retrieve}) chaque image qu'ensuite. Les \texttt{grab} s'enchaînant en quelques millisecondes, les trois prises sont \emph{quasi simultanées} --- \emph{indépendamment} de la cadence des caméras ---, un écart négligeable devant le mouvement des objets, immobiles au moment de la mesure. Une caméra qui échoue rend simplement une image \emph{absente}, que le pipeline gère sans s'interrompre.

\paragraph{Écarter les images périmées.} Les caméras USB \emph{tamponnent} quelques images ; lues telles quelles juste après un déplacement du bras, elles renverraient une vue \emph{périmée}. Avant chaque mesure, le pilote \emph{purge} donc ce tampon, de sorte que la reconstruction travaille sur l'état courant de la scène et non sur une image capturée en plein mouvement.

\dansLeProjet{Dans ce travail, les trois caméras USB --- réglées à $15$ images/s pour la paire stéréo et $10$ pour la caméra de poignet, précisément pour tenir la bande passante et éviter les décrochages --- sont réparties sur \emph{deux} concentrateurs (\emph{hubs}) : \texttt{cam\_0}/\texttt{cam\_1} d'un côté, \texttt{cam\_2} de l'autre. L'écart de saisie obtenu entre \texttt{cam\_0} et \texttt{cam\_1}, de quelques millisecondes, suffit pour une scène statique ; il ne conviendrait pas à un suivi d'objet en mouvement rapide, hors du cadre de ce travail.}

\section{Validation expérimentale}
\label{sec:perception-valid}

Cette section mesure la \emph{qualité géométrique} de la perception --- avec quelle précision la calibration situe les caméras et fonde la reconstruction. Elle ne traite pas de la réussite des prises : l'évaluation de la tâche complète, et la comparaison chiffrée des deux détecteurs sur le pick-and-place, relèvent du chapitre~\ref{ch:evaluation}.

\paragraph{Calibration intrinsèque.} L'erreur de reprojection finale (section~\ref{sec:calibration-cam}) reste \emph{sous-pixellique} pour les trois caméras --- $0{,}19$, $0{,}17$ et $0{,}14$ pixel pour \texttt{cam\_0}, \texttt{cam\_1} et \texttt{cam\_2} respectivement ---, chacune sur une quarantaine de vues du damier. Le modèle sténopé et sa distorsion décrivent donc fidèlement l'optique.

\paragraph{Calibration œil-main.} La qualité de la calibration extrinsèque se lit dans la \emph{cohérence} de la transformation constante recherchée : la pose du damier reconstruite à travers la chaîne devrait être identique d'une prise à l'autre, et l'écart résiduel mesure ce qui subsiste (table~\ref{tab:handeye-valid}). Trois lectures s'en dégagent. La caméra de poignet (\texttt{cam\_2}, \emph{eye-in-hand}) est de loin la plus précise --- $2{,}5$ mm et $0{,}4^{\circ}$ en moyenne ---, car elle observe le damier de près. Les caméras fixes, plus éloignées, sont plus lâches ($6{,}6$ mm pour \texttt{cam\_0}). Enfin \texttt{cam\_1} a bien ses \emph{intrinsèques} propres, mais sa \emph{pose} est déduite de \texttt{cam\_0} et de la transformation stéréo (section~\ref{sec:calibration-cam}) : son résidu cumule alors l'écart de \texttt{cam\_0} et celui du recalage stéréo, d'où une valeur plus élevée ($11{,}8$ mm). Ces résidus individuels, de l'ordre du demi-centimètre, restent inférieurs à la \emph{marge} d'environ $1$ cm que le planificateur laisse de \emph{chaque côté} de l'objet (chapitre~\ref{ch:planification}) --- d'autant que le raffinement final par la caméra de poignet corrige le reliquat avant la fermeture (chapitre~\ref{ch:controle}).

Reste la question que ces chiffres soulèvent : pourquoi calibrer \texttt{cam\_1} par \emph{déduction} plutôt qu'\emph{indépendamment} ? On l'a d'abord fait séparément, et c'est instructif. Chaque caméra affichait alors un bon résidu individuel (de l'ordre de $6$ mm), mais une \emph{cohérence stéréo} désastreuse : les deux caméras plaçaient un même point 3D à environ $21$ mm l'une de l'autre, soit un biais systématique d'environ $+40$ mm en $Y$ à la triangulation. La cause est géométrique --- deux calibrations indépendantes portent des erreurs \emph{sans corrélation} qui, sur une base orientée surtout selon $Y$, s'\emph{additionnent} au lieu de se compenser ---, alors que pour trianguler, seul l'\emph{accord} entre les deux caméras compte. La calibration \emph{stéréo conjointe} corrige exactement cela : elle mesure la transformation \emph{relative} $T_{\text{cam0}\leftarrow\text{cam1}}$ directement, puis en déduit \texttt{cam\_1}, de sorte qu'un biais de \texttt{cam\_0} devient \emph{commun} aux deux caméras et s'\emph{annule} à la triangulation. La cohérence du couple tombe ainsi à environ $2$ mm (biais résiduel ramené à environ $+14$ mm), au prix --- assumé --- du résidu individuel plus élevé de \texttt{cam\_1}. Autrement dit, le résidu individuel n'est pas la bonne mesure ici : c'est la cohérence \emph{relative}, seule utile à la triangulation, qui départage les deux méthodes.

\begin{table}[t]
  \centering
  \caption{Cohérence de la calibration œil-main : écart résiduel de la pose reconstruite sur les prises retenues (déviation moyenne et maximale), par caméra.}
  \label{tab:handeye-valid}
  \small
  \setlength{\tabcolsep}{4.5pt}
  \begin{tabular}{lcccc}
    \toprule
    Caméra & Montage & Poses (ret./capt.) & Transl. moy./max (mm) & Rot. moy./max ($^{\circ}$) \\
    \midrule
    \texttt{cam\_0} & \emph{eye-to-hand}           & $23/71$ & $6{,}6\ /\ 20{,}8$  & $1{,}4\ /\ 2{,}4$ \\
    \texttt{cam\_1} & \emph{eye-to-hand} (déduite) & $71/71$ & $11{,}8\ /\ 26{,}8$ & $2{,}9\ /\ 7{,}2$ \\
    \texttt{cam\_2} & \emph{eye-in-hand}           & $20/41$ & $2{,}5\ /\ 4{,}8$   & $0{,}4\ /\ 0{,}8$ \\
    \bottomrule
  \end{tabular}
\end{table}

\paragraph{Reconstruction 3D.} La précision de la triangulation peut se vérifier en comparant la position reconstruite d'un objet à sa position \emph{mesurée au pied à coulisse} depuis la base (outil \texttt{scripts/check\_perception.py}). Faute d'une campagne de mesures systématique, on s'en tient ici à un constat \emph{qualitatif}, conforme à l'analyse de la section~\ref{sec:reconstruction} : l'écart est dominé par la \emph{profondeur} --- l'axe caméra ---, les directions latérales étant plus fiables. La seule grandeur \emph{chiffrée} disponible est la \emph{cohérence stéréo} du couple, de l'ordre de $2$ mm (ci-dessus) ; les conséquences pratiques de l'incertitude 3D sur la réussite des prises sont examinées avec l'évaluation de la tâche (chapitre~\ref{ch:evaluation}).
